% -----------------------------------------------------------------------------
% Osamělá jednopísmenná spojky a předložky (a, i, k, o, s, u, v, z, …) na konci
% řádku: řádek se raději zalomí dříve a krátké slovo zůstane s následujícím slovem.
%
% LuaLaTeX: transform „oneletter.nobreak“ z balíčku babel (locale češtiny).
% Končí-li odstavec krátkým slovem, za ním už není slovní mezera, transform se
% neprojeví — odstavec může končit jedním písmenem.
%
% Dvoupísmenná slova tento mechanismus nepokrývá (v locale je jen „oneletter“).
% U častých tvarů (že, ve, se, …) doplňte v textu pevnou mezeru (~), případně
% použijte externí nástroj vlna na zdrojové soubory.
%
% pdfLaTeX: totéž z babelu nejde (transformy vyžadují LuaTeX); pro automatiku
% přejděte na LuaLaTeX (např. latexmk -lualatex) nebo doplňujte ~ ručně.
% -----------------------------------------------------------------------------
\IfFileExists{iftex.sty}{%
  \RequirePackage{iftex}%
}{%
  \PackageWarning{typografie-zlomy}{Chybí balíček iftex; podmíněná aktivace přeskočena.}%
}
\ifLuaTeX
  \AtBeginDocument{%
    \babelprovide[transforms=oneletter.nobreak]{czech}%
  }%
\fi
